\section{Hardware Mapping and Precision}
\label{sec:hardware-precision}

\subsection{From the adopted QUBO to an Ising Hamiltonian}

The adopted Q2 and Q3 formulations share the 104-variable
upper-triangular QUBO representation introduced in
Section~\ref{sec:q2-qubo}.  Let
\begin{equation}
 E_Q(\bm{x})=o+\sum_i Q_{ii}x_i+\sum_{i<j}Q_{ij}x_ix_j,
 \qquad x_i\in\{0,1\},
 \label{eq:hardware-qubo}
\end{equation}
where $s_i=2x_i-1$, or equivalently $x_i=(1+s_i)/2$.  Direct expansion
then gives the Ising form
\begin{equation}
 E_I(\bm{s})=K+\sum_i h_i s_i+\sum_{i<j}J_{ij}s_is_j,
 \label{eq:ising-hamiltonian}
\end{equation}
where
\begin{align}
 J_{ij}&=\frac{Q_{ij}}{4},\qquad
 h_i=\frac{Q_{ii}}{2}+\frac14\left(\sum_{j<i}Q_{ji}+\sum_{j>i}Q_{ij}\right),
 \label{eq:qubo-ising-coefficients}\\
 K&=o+\frac12\sum_iQ_{ii}+\frac14\sum_{i<j}Q_{ij}.
 \label{eq:qubo-ising-constant}
\end{align}
The constant $K$ does not affect the optimizer, but it must be retained
when comparing an Ising energy numerically with the original QUBO energy.
Kaiwu SDK~1.4.3 uses a reference-spin convention: it adds an auxiliary spin
$s_a$, writes $x_i=(1+s_is_a)/2$, and stores each pair twice in the returned
symmetric matrix.  The corresponding off-diagonal entries are
$-J_{ij}/2$, while the reference-spin entries are $-h_i/2$.  The
independently recomputed objective remains the one defined by
Equations~\eqref{eq:hardware-qubo}--\eqref{eq:qubo-ising-constant}.
Documenting this convention avoids offset and factor-of-two errors during
decoding.  Although QUBO--Ising equivalence is standard, implementation-
specific coefficient conventions still require an explicit audit
\cite{Lucas2014IsingNP,Glover2019QUBOTutorial}.

\subsection{Why direct eight-bit quantization is unsafe}

Fixed-coefficient hardware supports only a limited integer range.
Scaling a dense floating-point objective directly into that range can erase
small material-score differences while leaving large one-hot penalties
intact.  We therefore used deliberately artificial matrices in the
precheck only as stress tests; their quantized rankings are not treated as
Q2/Q3 evidence.  The tests illustrate two possible failures: the
quantized minimum can become infeasible, and a matrix can remain numerically
valid without preserving the intended material ranking.  This is
consistent with the general warning that penalty strength and coefficient
resolution must be assessed together
\cite{Ceselli2023GoodQUBOEncodings,Muecke2025OptimumPreservingQUBO,Ohno2026PrecisionErrorsQA}.

The adopted construction does not quantize the raw proxy scores
themselves.  Instead, within each fixed-$X$ group, it maps the capped,
normalized difference to the binary cost $c_k\in\{0,4\}$ in
Equation~\eqref{eq:q23-integer-cost}.  The known group winner receives zero,
and every other legal state receives four.  Combining these costs with the
two one-hot penalties ($\lambda=2$) and the offset ($o=4$) gives
Equation~\eqref{eq:q23-accepted-qubo}.  For legal states,
\begin{equation}
 E_Q(\bm{x}_{ab})=c_{ab},
 \qquad
 \arg\min_{a,b}E_Q(\bm{x}_{ab})=
 \arg\min_{a,b}\widehat V_{23}(a,b,X),
 \label{eq:group-winner-preservation}
\end{equation}
provided the group minimum is unique and every nonminimum gap is greater than
0.01.  Under NumPy \texttt{rint}, gaps below 0.01 and the exact half-step
0.01 both map to zero; see Equation~\eqref{eq:q23-integer-cost}.  The strict
condition was checked by exhaustive enumeration of all 2,704 legal states in
each adopted matrix.  The construction preserves
the group winner, not the complete raw-score geometry: non-winning values
collapse onto the same integer level, so their pairwise differences and
full ordering cannot be recovered from QUBO energies.  Cross-group selection
must therefore use the original-scale $\widehat V_{23}$ as specified in
Equation~\eqref{eq:q3-original-scale-comparison}.  Moreover, the level
boundary is defined using the already known group minimum $m_X$.  The adopted
two-level construction is therefore an oracle-assisted mapping and solver
test, not an independent method for discovering the proxy optimum.

\subsection{Coefficient and energy verification}

We use separate mathematical definitions for coefficient and energy
precision.  Let $(h_i^{\mathrm{ref}},J_{ij}^{\mathrm{ref}})$ be obtained by
directly applying Equations~\eqref{eq:qubo-ising-coefficients}--
\eqref{eq:qubo-ising-constant}, and let
$(h_i^{\mathrm{SDK}},J_{ij}^{\mathrm{SDK}})$ be reconstructed from the SDK
matrix after accounting for its reference-spin convention and duplicated
symmetric entries.  The coefficient residual is
\begin{equation}
 \epsilon_{\mathrm{coef}}=
 \max\left\{
   \max_i|h_i^{\mathrm{SDK}}-h_i^{\mathrm{ref}}|,
   \max_{i<j}|J_{ij}^{\mathrm{SDK}}-J_{ij}^{\mathrm{ref}}|
 \right\}.
 \label{eq:coefficient-residual}
\end{equation}
For a binary state $\bm x$, let $E_Q(\bm x)$ denote the full QUBO energy
including its offset, and let $E_{\mathrm{SDK}}^{\mathrm{corr}}(\bm x)$
denote the decoded SDK Ising energy with the conversion bias and QUBO
offset restored exactly once.  Over the audited state set
$\mathcal S_{\mathrm{audit}}$, the energy-chain residual is
\begin{equation}
 \epsilon_{\mathrm{energy}}=
 \max_{\bm x\in\mathcal S_{\mathrm{audit}}}
 \left|E_{\mathrm{SDK}}^{\mathrm{corr}}(\bm x)-E_Q(\bm x)\right|.
 \label{eq:energy-chain-residual}
\end{equation}
The constant is checked separately by
\begin{equation}
 \epsilon_{\mathrm{constant}}=|b+o-K_{\mathrm{ref}}|,
 \label{eq:constant-residual}
\end{equation}
where $b$ is the SDK conversion bias, $o$ is the QUBO offset, and
$K_{\mathrm{ref}}$ is the constant obtained from
Equation~\eqref{eq:qubo-ising-constant}.  The archived floating-point
controls satisfy $b+o=K_{\mathrm{ref}}$.  These residuals measure numerical
agreement in the mathematical QUBO--Ising conversion; they are not the
coefficient-quantization error or a measure of physical prediction
accuracy.

Expansion of the adopted integer QUBO gives entries in $\{-2,0,4\}$
before conversion.  For every adopted group, independent reconstruction
found
\begin{equation}
 \max_i|h_i|=102,
 \qquad \max_{i<j}|J_{ij}|=1,
 \qquad E_{\min}^{\mathrm{illegal}}-E_{\min}^{\mathrm{legal}}=2.
 \label{eq:formal-precision-audit}
\end{equation}
The coefficients fit within the signed eight-bit range without additional
rescaling.  For Q2 and each Q3 group, the exact formula-to-stored-matrix
energy residual over all 2,704 legal assignments was zero integer units.
A comparable Kaiwu~1.4.4 floating residual for the adopted 104-variable
integer matrices was not measured and remains NA.

For the separate floating-point controls, the measured
$(\epsilon_{\mathrm{coef}},\epsilon_{\mathrm{energy}})$ values were
$(2.13\times10^{-14},1.42\times10^{-12})$ for the 104-variable Q2
control and $(9.95\times10^{-14},4.57\times10^{-12})$ for the
208-variable Q3 control.  These are numerical round-off residuals from
the floating-point mapping audits; they must not be substituted for the
missing residual of the adopted 104-variable integer matrices.
The unified object/version/matrix/offset audit is given in
Table~\ref{tab:mapping-audit} in the appendix.

Kaiwu~1.4.4 was used only for the two floating mapping controls because no
solve licence was available.  The CPU sampling results used SDK~1.4.3.

The adopted workflow proceeds as follows: generate the versioned QUBO;
convert it using the documented SDK convention; verify the coefficient
range; solve with Kaiwu SDK~1.4.3 on the CPU; decode the reference spin;
check both one-hot constraints; and independently recompute $E_Q$ and the
original proxy score.  Each matrix was sampled with 30 seeds, ten reads per
seed, 300 iterations, and $\alpha=0.99$.  We accepted no solver-reported
energy as a material result without this independent decoding.
Sections~\ref{sec:q2-qubo} and~\ref{sec:q3-qubo} give the decoded candidates
and sample statistics.

\subsection{Projection, quantization, and ranking controls}

Three controls show why coefficient range alone is insufficient.  First, the
107-variable second-order Q3 projection achieved Spearman 0.978891 over all
8,112 candidates but changed the optimum from Na--Nb--O3
($\widehat V_{23}=0.540023$) to Rb--Ta--O3 (0.676651), had zero Top-5 overlap,
and produced $\max|h|=254$ after scaling.  Its global fit therefore did not
preserve the decision-relevant low tail.

Second, uniform full-pool quantization retained the raw-score
minimizer inside the minimum integer-energy tier, but did not make it the
unique selected candidate.  Let
$\mathcal T_{\min}=\{z:E_{\mathrm{int}}(z)=\min_uE_{\mathrm{int}}(u)\}$
and $z^*=\arg\min_z\widehat V_{23}(z)$.  We distinguish the indicators
$\mathbb I[z^*\in\mathcal T_{\min}]$ and
$\mathbb I[z^*=z_{\mathrm{ID}}]$, where the deterministic lexicographic
tie-break is
$z_{\mathrm{ID}}=\arg\min_z(E_{\mathrm{int}}(z),\mathrm{ID}(z))$.
For the uniform qubo_x127/4 mapping, the first indicator is one
but the second is zero: the minimum tier contains six Q2 candidates
and 21 Q3 candidates.  The qubo_xfloor Q3 tier contains 24 candidates.
The reported score loss is
\begin{equation}
 \Delta_{\mathrm{tie}}\widehat V
 =\widehat V_{23}(z_{\mathrm{ID}})-\widehat V_{23}(z^*),
 \label{eq:tie-break-score-loss}
\end{equation}
which equals $0.647403$ for Q2 and $0.645059$ for Q3.  It is a
prediction-score loss caused by breaking an integer-energy tie by ID,
not an energy difference and not a QUBO--Ising mapping residual.
The pairwise order-retention rate among candidate pairs assigned
different integer energies was $1.000$; overall pairwise retention was
lower (about $0.918$ for Q2 and $0.924$--$0.926$ for Q3) because ties
collapse distinctions.  Thus
``the optimum is in the lowest tier'' and ``the quantized model returns the
unique original optimum'' are different claims.

Third, the 18-point $K\times\lambda$ scan exposed a feasibility--range
trade-off.  For $\lambda=2,4,8$ the lowest illegal-minus-legal gaps were
$-2,0,4$; at $K=52,\lambda=8$, however, $\max|h|=252$ exceeded the signed
eight-bit range.  The historical stage-2 record of 0/300 legal samples and the later
8/300 record were traced to a sample-counting defect, not to different
QUBO matrices or solver behaviour.  On the same 300 returned samples
(seed 12345, identical matrix and parameters), the corrected counter
identified 8 legal samples and 8 proxy-minimum hits, whereas the legacy
counter returned zero because it subtracted the left-block width from an
already absolute column index.  The stage-2 and stage-3 entry points
produced identical matrices and variable mappings.  The original historical
0/300 run log is not archived, so that historical result remains
\textbf{UNVERIFIED}; it is not treated as a solver-performance estimate.
A separate explicit-seed reconciliation (O11-R03) recorded 10 legal
samples and 10 hits in 300 reads, with at least one hit in 9 of 30
ten-read batches.  Sample-level counts and batch-level hit rates are
reported separately.

\subsection{Final-matrix SPQC-550 observations}

In addition to the historical 17- and 30-variable short-list tasks, four
final 104-variable matrices were submitted to SPQC-550.  Their task-level
returns and timing fields are in Appendix Table~\ref{tab:spqc-final-results}.

Three of four tasks returned at least one group proxy argmin among unequal
returns of 10, 10, 7, and 5 samples.  This count is not a hardware success
rate.  The NO2 Q3 task is
retained as a zero-legal-sample counterexample rather than discarded.  The
four submitted matrix hashes and complete task IDs are listed in the
appendix.  The quota ledger records seven debits from a balance change of 31
to 24; six task receipts are archived.  The additional \texttt{probe4\_01}
entry lacks an original receipt, so its task fields remain NA.  Platform queue time, full execution time, and
monetary billing remain NA.

\subsection{Evidence boundary}

The solver statistics in Sections~\ref{sec:q2-qubo} and
\ref{sec:q3-qubo} are Kaiwu SDK~1.4.3 CPU results.  The four SPQC
submissions returned the samples listed in Table~\ref{tab:spqc-final-results},
but their small and unequal sample counts do not support estimates of
stability or speedup.  Blind-pool and projection experiments are also
local comparisons.  Every reported candidate remains a proxy argmin,
not a physically measured optimum.
